% ============================================================
% P16 — Retur dan Refund
% ============================================================
\flowmeta{P16}{Retur dan Refund}{Pelanggan}

\begin{flowsection}{Tujuan}
Mengajukan pengembalian barang (retur) dan permintaan pengembalian dana (refund)
atas pesanan yang telah selesai.
\end{flowsection}

\begin{flowsection}{Prasyarat}
Sudah login dan memiliki pesanan yang memenuhi syarat retur (mis. dalam masa
garansi tukar ukuran 14 hari).
\end{flowsection}

\begin{flowsection}{Langkah Mengajukan Retur}
\begin{enumerate}
  \item Buka \texttt{/returns/request}.
  \item Pilih pesanan dan produk yang ingin diretur.
  \item Isi alasan retur dan jumlah barang.
  \item Kirim pengajuan; tunggu peninjauan admin.
\end{enumerate}
\end{flowsection}

\shot{gambar/pelanggan/p16-1.png}{Formulir pengajuan retur \& refund}{fig:p16-1}

\begin{flowsection}{Hasil yang Diharapkan}
Pengajuan tercatat dengan status menunggu; pelanggan dapat memantau perkembangannya
melalui detail pesanan.
\end{flowsection}

\begin{flowsection}{Penanganan Error dan Kasus Khusus}
\begin{itemize}
  \item Pesanan tidak memenuhi syarat retur $\rightarrow$ pengajuan ditolak.
  \item Belum ada pesanan $\rightarrow$ tidak ada yang dapat diretur.
\end{itemize}
\end{flowsection}

\begin{flowsection}{Kaitan Modul}
FE \texttt{ReturnRequestPage.jsx}, \texttt{ReturnDetailPage.jsx}. BE
\texttt{ReturnController}. Rute \texttt{/api/returns}.
\end{flowsection}

\docver{v1.0}{Sep 2026}
